
% 中英文摘要和关键字
% 中文摘要
\renewcommand{\abstractname}{摘要} % 将 abstract 标题改为“摘要”
\begin{abstract}
集成感知与通信（ISAC）已被公认为下一代B5G及6G移动网络的核心支柱技术。本文旨在为ISAC提供一份全面的技术蓝图。文章首先通过回顾雷达与通信的融合历史，明确了ISAC在提高频谱效率、硬件效率及性能增益方面的优势。随后，本文系统性地分析了通感一体化在信息论、物理层及跨层设计中的基础性能权衡。在技术实现层面，重点探讨了面向双功能系统的联合波形设计与先进信号处理算法。文章进一步阐述了“感知网络”的演进范式，论证了感知与通信如何通过相互辅助实现系统级性能提升。最后，本文展望了ISAC与边缘智能等技术的集成前景，为未来全感知无线网络的设计提供了重要参考。


\end{abstract}

% 中文关键词
\noindent \textbf{关键词：} 
集成感知与通信（ISAC）、6G、性能权衡（Performance Tradeoff）、波形设计、感知网络（Perceptive Network）
\addcontentsline{toc}{section}{摘要} % 将摘要添加到目录中

\; 

\renewcommand{\abstractname}{Abstract} % 恢复为默认的“Abstract”


\begin{abstract}
Integrated Sensing and Communication (ISAC) has been widely recognized as a key enabling technology and a fundamental pillar of next-generation Beyond-5G (B5G) and 6G mobile networks. This paper aims to provide a comprehensive technical blueprint for ISAC. It begins by reviewing the historical evolution of radar–communication convergence, highlighting the advantages of ISAC in terms of improved spectrum efficiency, hardware efficiency, and performance gains. The paper then systematically analyzes the fundamental performance trade-offs of ISAC from the perspectives of information theory, physical-layer design, and cross-layer optimization. At the implementation level, particular emphasis is placed on joint waveform design and advanced signal processing algorithms for dual-functional sensing and communication systems. Furthermore, the paper elaborates on the emerging paradigm of perceptive networks, demonstrating how sensing and communication can mutually assist each other to achieve system-level performance enhancement. Finally, it discusses the integration of ISAC with emerging technologies such as edge intelligence, reconfigurable intelligent surfaces (RIS), providing valuable insights for the design of future fully perceptive wireless networks.


\end{abstract}

% 英文关键词
\noindent \textbf{Keywords:}
Integrated Sensing and Communication (ISAC), 6G, Performance Tradeoff, Waveform Design, Perceptive Network

\addcontentsline{toc}{section}{Abstract} % 将摘要添加到目录中
